\chapter{From Bellman to Preference Learning: The Grand Arc}
\label{ch:grand-arc}

This chapter tells a single story.
It begins in 1957 with Richard Bellman's dynamic programming and ends in 2025 with DeepSeek's GRPO---a method in which the Bellman equation appears nowhere.
Along the way, we trace nine stages of evolution and identify the forces that drove each transition.

The thesis in one sentence:

\begin{keyidea}[The Grand Arc]
Classical DP solves for the global optimum under a known model.
Modern RL learns a good-enough, stable, preference-aligned policy under limited sampling, noise, bias, and distributional uncertainty.
These two objectives differ in nature: the former is god's-eye analytical optimization; the latter is statistical preference learning.
\end{keyidea}

% ===========================================================
\section{Stage 1: The Bellman Equation---Correct, but Ideal}
\label{sec:stage1-dp}
% ===========================================================

In 1957, Richard Bellman introduced dynamic programming and the principle of optimality.
In a deterministic, fully known environment with transition model $P(s'|s,a)$ and reward $r(s,a)$, the Bellman equation is both necessary and sufficient.

For a given policy $\pi$:
\begin{equation}
V^{\pi}(s) = \mathbb{E}_{a \sim \pi,\, s' \sim P}\bigl[r(s,a) + \gamma\, V^{\pi}(s')\bigr].
\end{equation}
For optimal control:
\begin{equation}
V^{\ast}(s) = \max_a\, \mathbb{E}_{s' \sim P}\bigl[r(s,a) + \gamma\, V^{\ast}(s')\bigr].
\end{equation}
Or in Q-function form:
\begin{equation}
\label{eq:bellman-q}
Q^{\ast}(s,a) = \mathbb{E}_{s' \sim P}\bigl[r(s,a) + \gamma\, \max_{a'} Q^{\ast}(s',a')\bigr].
\end{equation}

These equations are mathematically rigorous.
The Bellman equation \emph{is} the entire optimization---solving it yields the globally optimal policy.
The problem is not that the Bellman principle is wrong, but that it describes an \textbf{ideal fixed point}, while the training process is not a fixed point.

% ===========================================================
\section{Stage 2: Stochastic DP and MDPs (1960)}
\label{sec:stage2-mdp}
% ===========================================================

Ronald Howard extended Bellman's framework to stochastic environments (Markov Decision Processes) and introduced Policy Iteration and Value Iteration.
The essence did not change---still finding the Bellman equation's fixed point under a known model---but transition probability uncertainty was now accommodated.

The Bellman evaluation equation becomes:
\begin{equation}
V^{\pi}(s) = \sum_a \pi(a|s) \sum_{s'} P(s'|s,a)\bigl[r(s,a) + \gamma\, V^{\pi}(s')\bigr].
\end{equation}

This operator is fundamentally \emph{averaging}: it sums over actions weighted by the policy and over next states weighted by transition probabilities.
Averaging is smooth, stable, and easy to prove contraction for.
This stability is why V-functions were developed first.

% ===========================================================
\section{Stage 3: TD Learning---From Known Models to Sampling (1988)}
\label{sec:stage3-td}
% ===========================================================

Sutton's Temporal Difference (TD) learning was the first major departure from the god's-eye view.
The key move: replace exact expectations with sampled estimates.

\begin{equation}
V(s) \leftarrow V(s) + \alpha\bigl[r + \gamma\, V(s') - V(s)\bigr].
\end{equation}

No model is needed---only the ability to observe transitions $(s, a, r, s')$.
V-functions were the natural starting point because they average over action outcomes, yielding low-variance, stable updates.

The cost: V cannot directly tell you which action to take.
To extract a policy from $V$, you need the transition model $P(s'|s,a)$---precisely what you no longer have.

% ===========================================================
\section{Stage 4: Q-Learning---Making Action Optimization Explicit (1989)}
\label{sec:stage4-q}
% ===========================================================

Watkins' Q-learning solved the action-selection problem by shifting the learning target from $V(s)$ to $Q(s,a)$:
\begin{equation}
Q(s,a) \leftarrow r + \gamma\,\max_{a'} Q(s',a').
\end{equation}

Policy improvement becomes trivial: $a^{\ast} = \arg\max_a Q(s,a)$.
No environment model needed.
Q-learning is also \textbf{off-policy}---the update assumes the future will choose optimally regardless of what the behavior policy actually did.
This makes it data-efficient: a high-value signal can propagate rapidly.

\subsection{The Cost of $\max$}

But the $\max$ operator is an unbridled horse.
During training, Q is a noisy estimate: $\hat{Q} = Q^{\ast} + \epsilon$.
Therefore:
\begin{equation}
\max_a \hat{Q}(s,a) = \max_a \bigl[Q^{\ast}(s,a) + \epsilon(s,a)\bigr].
\end{equation}
The $\max$ preferentially selects the action with the largest positive noise.
This creates a feedback loop:
\begin{enumerate}
    \item An action is overestimated due to noise.
    \item $\max$ selects it for the Bellman target.
    \item The inflated target is written back into the Q-network.
    \item Future updates propagate the overestimation further.
\end{enumerate}

\begin{keyidea}[The Core Tension]
\textbf{Averaging} (V-style) is stable but conservative---it does not amplify the best action.
\textbf{Max} (Q-style) is data-efficient but aggressive---it amplifies estimation noise.
The entire history of RL oscillates between these two poles.
\end{keyidea}

\begin{table}[ht]
\centering
\caption{The averaging--max tension.}
\label{tab:avg-vs-max}
\begin{tabular}{lll}
\toprule
\textbf{Approach} & \textbf{Strengths} & \textbf{Weaknesses} \\
\midrule
Averaging / expectation & Stable, smooth, low variance & Slow learning, weak optimization \\
Max / greedy & Data-efficient, fast propagation & Overestimation, instability \\
\bottomrule
\end{tabular}
\end{table}

% ===========================================================
\section{Stage 5: DQN---Taming Qmax with Deep Networks (2013--2016)}
\label{sec:stage5-dqn}
% ===========================================================

DQN did not invent Q-learning.
Its true significance is the first reasonably stable combination of deep neural networks with off-policy Q-learning.
The Bellman target remained unchanged:
$y = r + \gamma\,\max_{a'} Q_{\bar\theta}(s',a')$.
DQN tamed the $\max$ operator through engineering:

\begin{enumerate}
    \item \textbf{CNN encoder}: map high-dimensional pixel states to features.
    \item \textbf{Replay buffer}: break temporal correlation, improve data reuse.
    \item \textbf{Target network} $\bar\theta$: prevent the target from chasing itself, reducing bootstrap instability.
    \item \textbf{Off-policy learning}: reuse past data to train the current value function.
\end{enumerate}

\subsection{Subsequent Repairs}

DQN exposed the $\max$ operator's weaknesses at scale.
Two key follow-ups:

\textbf{Double DQN} decouples action selection from evaluation:
\begin{equation}
a^{\ast} = \arg\max_a Q_\theta(s',a), \qquad
y = r + \gamma\, Q_{\bar\theta}(s', a^{\ast}).
\end{equation}
The online network selects; the target network evaluates.
This reduces the noise amplification from using the same noisy Q for both roles.

\textbf{Dueling DQN} decomposes the Q-function structurally:
\begin{equation}
Q(s,a) = V(s) + A(s,a),
\end{equation}
where $V(s)$ captures overall state value and $A(s,a)$ captures the relative advantage of each action.
By isolating the state-value component, estimation becomes more stable.

\begin{keyidea}[DQN's Legacy]
The DQN family retains Qmax's aggressive optimization but puts reins on it.
The Bellman equation is still the primary objective---but it now requires heavy engineering stabilization.
\end{keyidea}

% ===========================================================
\section{Stage 6: Policy Gradient---The Philosophical Watershed}
\label{sec:stage6-pg}
% ===========================================================

Q-learning is fundamentally \textbf{value-centric}: it posits an objectively optimal $Q^{\ast}(s,a)$ and approaches it through Bellman consistency.
Policy Gradient (Williams 1992; Sutton et al.\ 1999) is entirely different.
It directly optimizes the policy $\pi_\theta(a|s)$.

The objective:
\begin{equation}
J(\theta) = \mathbb{E}_{\tau \sim \pi_\theta}\bigl[R(\tau)\bigr].
\end{equation}
The gradient (via the likelihood ratio trick):
\begin{equation}
\label{eq:pg}
\nabla_\theta J(\theta) = \mathbb{E}\Bigl[\nabla_\theta \log \pi_\theta(a|s)\, A(s,a)\Bigr].
\end{equation}

The meaning is intuitive: if an action's advantage is positive, increase its probability; if negative, decrease it.
This is not hard $\max$---it is \textbf{probability reshaping}.

\subsection{Two Layers of Probability}

The statistical essence of policy gradient can be understood through two layers:
\begin{enumerate}
    \item \textbf{Rollout probability}: which actions get sampled depends on $\pi_\theta(a|s)$.
    \item \textbf{Update weighting}: sampled actions are re-weighted by advantage $A(s,a)$.
\end{enumerate}
Policy gradient does not simply copy sampled behaviors.
Within the actual sampling distribution, it re-weights behaviors by reward as preferences---closely mirroring statistical estimation.

\subsection{A Change of Objective}

The shift from Q-learning to policy gradient is not merely algorithmic---the optimization objective changes.

Q-learning's Bellman loss is a \textbf{consistency loss}:
\begin{equation}
\mathcal{L}_Q = \bigl(Q(s,a) - r - \gamma\,\max_{a'} Q(s',a')\bigr)^2.
\end{equation}
It pursues value-function self-consistency---fitting ``true'' values.

Policy gradient's loss is \textbf{behavior distribution shaping}:
\begin{equation}
\mathcal{L}_\pi = -\log \pi(a|s)\, A(s,a).
\end{equation}
It pursues increasing good-behavior probabilities and decreasing bad ones.

\begin{keyidea}[The Watershed]
Q-learning says: ``The action with the highest estimated value should be selected.''
Policy gradient says: ``Behaviors that perform statistically better should occur more frequently.''
This is the turn from god's-eye value solving to sampling-based behavioral preference shaping.
\end{keyidea}

\subsection{Policy Gradient as Weighted Maximum Likelihood}

The connection to statistics is precise.
Ordinary maximum likelihood (behavior cloning) maximizes $\sum_i \log \pi_\theta(a_i|s_i)$---every observed action is reinforced equally.
Policy gradient maximizes a \emph{weighted} log-likelihood:
\begin{equation}
\sum_i A(s_i, a_i)\,\log \pi_\theta(a_i|s_i).
\end{equation}
MLE models behavioral frequency; policy gradient models behavioral frequency with preference weights.

% ===========================================================
\section{Stage 7: Actor-Critic and the PPO/SAC Era (2015--2018)}
\label{sec:stage7-ac}
% ===========================================================

Policy gradient needs advantage estimates.
Actor-Critic methods maintain both a policy (actor) and a value function (critic):

\begin{table}[ht]
\centering
\caption{Critic variants in Actor-Critic architectures.}
\label{tab:critic-types}
\begin{tabular}{lll}
\toprule
\textbf{Critic Type} & \textbf{What It Learns} & \textbf{Typical Algorithms} \\
\midrule
V-critic    & $V(s)$    & A2C, A3C, PPO \\
Q-critic    & $Q(s,a)$  & DDPG, TD3, SAC \\
\bottomrule
\end{tabular}
\end{table}

PPO-style methods use a V-critic to estimate advantage via TD residuals or GAE:
\begin{equation}
\hat{A}_t \approx r_t + \gamma\, V(s_{t+1}) - V(s_t).
\end{equation}
The critic here is a statistical control variate---not an attempt to find a globally optimal $Q^{\ast}$.

\subsection{The Key Turn: Policy Becomes Primary}

TRPO, PPO, and SAC marked a fundamental shift.
The optimization objective is no longer Bellman consistency but:
\begin{equation}
\max_\theta\; \mathbb{E}_{\pi_\theta}\bigl[R(\tau)\bigr].
\end{equation}
The Bellman equation is \emph{demoted to an auxiliary role}---the critic uses it to estimate advantage, but it is no longer the primary optimization formula.

\subsection{Greedy with Brakes}

Modern methods can be understood as \textbf{greedy optimization with brakes}:

\begin{itemize}
    \item \textbf{Why greedy?}  Data is scarce.  Without greedy amplification, high-reward signals propagate too slowly.
    \item \textbf{Why brakes?}  Greedy amplifies noise, bias, and function approximation error.
\end{itemize}

\begin{table}[ht]
\centering
\caption{How modern algorithms balance greediness and stability.}
\label{tab:brakes}
\begin{tabular}{ll}
\toprule
\textbf{Algorithm / Technique} & \textbf{Essential Function} \\
\midrule
PPO clipping              & Limit overly aggressive policy updates \\
TRPO / KL constraint      & Trust region prevents policy jumps \\
Entropy regularization    & Prevent premature deterministic collapse \\
TD3 clipped double Q      & Conservative estimates suppress overestimation \\
SAC                       & Soft greedy + entropy + Q-critic \\
Target network            & Reduce bootstrap target drift \\
Double Q                  & Reduce $\max$ overestimation \\
\bottomrule
\end{tabular}
\end{table}

The shared principle: \emph{do not abandon optimization, but prevent the optimizer from exploiting estimation noise.}

\subsection{V-Only as a Relaxation of Qmax}

PPO replaces hard $\arg\max$ with soft advantage-weighted updates:
\begin{align}
\text{Q-learning:} \quad & a^{\ast} = \arg\max_a Q(s,a), \\
\text{PPO:} \quad & \nabla_\theta \log \pi_\theta(a|s)\, A(s,a).
\end{align}
The former is hard optimization; the latter is soft preference adjustment.
PPO is more stable but not ``more globally optimal.''
What it pursues is: stably pushing the policy toward better directions under limited sampling.

% ===========================================================
\section{Stage 8: RLHF---Preferences Replace Engineered Rewards (2017--2022)}
\label{sec:stage8-rlhf}
% ===========================================================

Christiano et al.\ (2017) proposed learning a reward model from human preference comparisons.
Ouyang et al.\ (2022, InstructGPT) applied it at scale to LLMs.

The shift: reward is no longer a hand-designed function but is \emph{sampled from human preferences}---noisy, contextual, local, sometimes contradictory.
Forcing a globally consistent Bellman-optimal $Q^{\ast}$ onto such preferences may itself be a fiction.

However, RLHF's policy training still used PPO internally.
PPO's critic still relied on the Bellman equation for value estimation.
Preferences changed the \emph{source} of reward, but the Bellman equation had not yet fully exited the optimization mechanism.

% ===========================================================
\section{Stage 9: DPO and GRPO---The Bellman Equation Disappears (2023--2025)}
\label{sec:stage9-dpo}
% ===========================================================

This is the most dramatic rupture in the entire arc.

\textbf{DPO} (Rafailov et al., 2023) directly optimizes the policy on preference pairs, eliminating the reward model and the RL training loop entirely.
No critic, no value network, no Bellman bootstrap.

\textbf{GRPO} (DeepSeek, 2025) goes further: for each prompt, generate a group of responses and use within-group relative ranking as the baseline to replace the critic.
It is pure policy gradient plus statistical preference re-weighting.

\begin{keyidea}[The Complete Break]
In DPO/GRPO there is:
\begin{itemize}
    \item No value network ($V$ or $Q$).
    \item No Bellman consistency loss.
    \item No bootstrap target.
\end{itemize}
The optimization objective is pure preference distribution shaping.
The Bellman equation is completely gone.
\end{keyidea}

This is not regression---it is the inevitable endpoint of a trend.
In a world where human preferences are inherently noisy, contextual, and ad-hoc, maintaining a globally consistent Bellman fixed point is neither necessary nor natural.
Statistical preference optimization is the correct response to this reality.

% ===========================================================
\section{The Diminishing Role of the Bellman Equation}
\label{sec:bellman-role}
% ===========================================================

The entire nine-stage evolution can be summarized by tracking a single variable: \emph{how central is the Bellman equation to the optimization?}

\begin{table}[ht]
\centering
\caption{The Bellman equation's role at each stage.}
\label{tab:bellman-role}
\begin{tabular}{lll}
\toprule
\textbf{Stage} & \textbf{Era} & \textbf{Bellman's Role} \\
\midrule
1--4. DP / Q-learning       & 1957--1989 & The entire optimization \\
5. DQN family               & 2013--2016 & The entire optimization (+ stability tricks) \\
6. Early Actor-Critic       & $\sim$2000 & Primary objective (policy is a means) \\
7. PPO / SAC                & 2015--2018 & Auxiliary tool (critic uses it for advantage) \\
8. RLHF                     & 2017--2022 & Residual (inside PPO's critic) \\
9. DPO / GRPO               & 2023--2025 & Completely absent \\
\bottomrule
\end{tabular}
\end{table}

% ===========================================================
\section{The $\arg\max$ Distinction: Conclusion vs.\ Process}
\label{sec:argmax-distinction}
% ===========================================================

A subtle but critical point underlies the entire arc.

In Bellman optimality, $\arg\max$ is a \textbf{conclusion}:
\emph{if} you already know the true $Q^{\ast}$, \emph{then} choosing $\arg\max Q^{\ast}$ is correct.

During training, $\arg\max$ becomes the \textbf{process} at every step:
even though $Q$ is noisy, you forcefully select the current maximum.

These are completely different operations.
An analogy: Bellman optimality's $\arg\max$ is knowing that daily exercise is good \emph{when you are already fit}.
Hard greedy during training is attempting maximum-intensity exercise every day \emph{when your body is in poor condition}---overshooting, injury, collapse.

Modern RL makes the optimization ``softer'':
\begin{itemize}
    \item PPO uses clipping to prevent the policy from changing too much in one step.
    \item SAC uses entropy to prevent premature collapse.
    \item Double Q / TD3 uses conservative estimates to suppress overestimation.
    \item RLHF uses KL regularization to stay close to the reference policy.
\end{itemize}

None of these deny optimality.
They acknowledge that under limited data, you cannot treat theoretical optimality conditions as hard execution rules at every training step.

% ===========================================================
\section{From God's-Eye View to Statistical View}
\label{sec:god-to-stat}
% ===========================================================

Classical DP adopts a god's-eye view: there exists a true model, true transition probabilities, a true reward function, and true optimal $Q^{\ast}$, $V^{\ast}$.
The algorithm's task is to compute these true values.

The statistical view is different.
There is no directly observable truth---only samples, assumptions, likelihoods, and estimators.
Maximum likelihood estimation does not seek cosmic truth; it asks:
\emph{what parameters best explain the current observed data?}

As RL has evolved, it has increasingly adopted this statistical worldview:
\begin{itemize}
    \item Policy gradient, RLHF, and GRPO do not insist on a globally consistent $Q^{\ast}$.
    \item They care about: which behaviors brought better rewards in sampling statistics, and therefore should occur more frequently.
\end{itemize}

This is the shift from ``finding true values'' to ``shaping behavioral distributions.''

\begin{keyidea}[Two Worldviews]
DP is exact optimization under a known model---\textbf{analytical optimal control}.\\
Modern RL is recursive preference optimization under noise, bias, and distribution drift---\textbf{statistical preference learning}.
\end{keyidea}

% ===========================================================
\section{The Six Phases at a Glance}
\label{sec:six-phases}
% ===========================================================

The nine stages cluster naturally into six conceptual phases.

\subsection*{Phase 1: DP / V-Function}
\emph{Goal}: stable evaluation and globally optimal solving under known models.\\
\emph{Character}: exact expectations, clean Bellman fixed point, stable V-function.\\
\emph{Limitation}: action optimization is not direct.

\subsection*{Phase 2: Q-Learning / Qmax}
\emph{Goal}: explicit action optimization without a model.\\
\emph{Character}: high data efficiency, off-policy, fast greedy propagation.\\
\emph{Limitation}: severe overestimation, instability.

\subsection*{Phase 3: DQN Family}
\emph{Goal}: scale Q-learning with deep networks.\\
\emph{Character}: still value-centric, still Bellman-centered, but with heavy engineering stabilization.\\
\emph{Limitation}: $\max$ bias persists despite repairs.

\subsection*{Phase 4: Policy Gradient / Actor-Critic}
\emph{Goal}: directly optimize the policy distribution.\\
\emph{Character}: shift from value-centric to behavior-centric; no longer insisting on perfect Q.\\
\emph{Limitation}: high variance, lower data efficiency.

\subsection*{Phase 5: PPO / SAC / TD3}
\emph{Goal}: fuse value and policy, balance greediness and stability.\\
\emph{Character}: constrained/regularized policy updates; the Bellman equation becomes auxiliary.\\
\emph{Limitation}: still requires careful hyperparameter tuning.

\subsection*{Phase 6: RLHF / DPO / GRPO}
\emph{Goal}: align with human preferences.\\
\emph{Character}: reward from preferences (not engineering); algorithms resemble statistical preference learning.\\
\emph{Essence}: the goal is not absolute optimality, but placing the behavioral distribution within the zone of human preferences.

% ===========================================================
\section{Chapter Summary}
\label{sec:ch2-summary}
% ===========================================================

Ten points compress the grand arc:

\begin{enumerate}
    \item DP is exact optimization under known models; RL is sampling-based optimization under unknown models.
    \item The Bellman equation describes an ideal fixed point; the training process is noisy nonlinear recursive iteration.
    \item V is stable but conservative; Q enables explicit action optimization but more easily amplifies errors.
    \item $\max$ is data-efficient because it rapidly propagates high values---but it also preferentially selects positive noise.
    \item Modern RL's core is not pursuing global optimality, but learning stable, useful, statistically superior policies under limited data.
    \item Policy Gradient is the watershed: from finding the optimal value function to directly shaping behavior distributions.
    \item Actor-Critic fuses value estimation and policy optimization---not simply ``using Q'' or ``not using Q.''
    \item PPO / SAC / TD3 are greedy optimization with brakes.
    \item RLHF / GRPO demonstrate that real-world preferences are inherently ad-hoc and statistical---not suited for a globally Bellman-consistent value function.
    \item Modern RL more closely resembles statistical preference learning than classical analytical optimal control.
\end{enumerate}

The entire thread in one powerful statement:

\begin{keyidea}[The Essence]
DP pursues the global optimal solution.
RL pursues evolving good-enough behavioral preferences from limited, noisy, biased data.
The field's history oscillates between ``too much averaging---can't learn'' and ``too much greediness---collapses.''
The answer is not choosing one side, but:
use greediness to amplify effective signals; use statistical constraints and regularization to prevent the system from being carried away by noise.
\end{keyidea}
